\section{RELATED WORK}\label{sec:Related_Work}

\subsection{Federated Learning: Heterogeneity and Communication Constraints}

The approach of federated learning~\cite{mcmahan2017communication} represents a compelling solution for distributed machine learning that maintains data privacy. 
This approach, anchored by the fundamental FedAvg algorithm, facilitates collaborative model training without requiring consolidation of raw data, 
making it particularly suitable for sensitive applications like satellite telemetry. 
However, the effectiveness of federated learning in real-world scenarios is significantly challenged by two interrelated problems: data heterogeneity and communication constraints.

Data heterogeneity across clients has been extensively studied, with numerous approaches proposed to improve robustness. 
To tackle non-independent and identically distributed (non-IID) data distributions, FedProx~\cite{li2020federated} incorporates a proximal term. 
FedDyn~\cite{jin2023feddyn} employs dynamic regularization to handle heterogeneous objectives. 
FedAdam and other adaptive optimization methods~\cite{reddi2020adaptive} provide server-side momentum for improved convergence.

Communication efficiency is another critical bottleneck in federated learning deployments. 
Gradient compression techniques~\cite{lin2017deep}, quantization methods, and structured updates have been proposed to reduce communication overhead. 
However, most existing works assume clients can participate in multiple federated rounds. 
In contrast, satellite TT\&C systems face extreme communication constraints where some centers may only participate in a single round.

\subsection{Continual Learning and Knowledge Preservation in Federated Settings}

To address catastrophic forgetting, various continual learning methods have been developed. 
The class-incremental problem in federated learning—where new fault classes emerge over time—remains an open challenge in the learning community~\cite{de2021continual}. 
Several continual learning approaches have been proposed to mitigate catastrophic forgetting. 
Yoon et al.~\cite{yoon2021federated} divide model parameters into baseline and task-adaptive components to balance knowledge acquisition and retention. 
Similarly, Dong et al.~\cite{de2021continual} decompose gradients into new and old components for selective updates.

A number of regularization-based methods have been developed to maintain model performance on previous tasks. 
Knowledge distillation represents a valuable approach for maintaining model knowledge throughout the learning process. 
Despite its effectiveness in centralized settings, knowledge distillation has received limited attention in federated continual learning contexts.
Through the guidance of a teacher model during student model training with new data, knowledge distillation helps reduce performance deterioration on previously learned fault classes. 
In the federated learning context, knowledge distillation has been applied to preserve model knowledge and improve collaborative training efficiency.

In federated settings, teacher-student frameworks enable knowledge transfer while maintaining privacy. 
By using a teacher model to guide student model training on new data, knowledge distillation can mitigate performance degradation on previously learned fault classes. 
This approach aligns well with the privacy constraints of satellite systems, as it does not require sharing raw gradients or private data.

However, existing continual learning and knowledge distillation methods typically require sharing either raw gradients, specific parameter updates, or extensive model interactions. 
Recent studies on Deep Leakage from Gradients (DLG) demonstrate that original training data can be reconstructed solely from shared gradients~\cite{zhu2019deep, Geiping2020}. 
Thus, in TT\&C applications where data sovereignty is strictly enforced, such information sharing is prohibited. 
Most existing works that address both heterogeneity and privacy concerns~\cite{huang2022learn, jiang2021fedspeech} typically assume irrelevant knowledge across clients, 
whereas TT\&C centers require comprehensive fault knowledge sharing to maximize collective diagnostic capability.

\subsection{Fault Diagnosis in Satellite Telemetry Systems}

In aerospace engineering, significant research efforts have focused on data-driven approaches for fault diagnosis~\cite{ganesan2021fault, xiao2020deep}. 
Traditional approaches include statistical methods and domain-specific feature engineering. 
Contemporary developments leverage deep learning approaches including convolutional neural networks and time series transformers to achieve automated identification of fault patterns. 
The Dual-aspect Time Series Transformer (DTST) method~\cite{xu2024dtst} represents SOTA performance in space system fault diagnosis. 
However, these centralized approaches typically require aggregating sensitive telemetry data at a single location, which contradicts the data privacy requirements in satellite operations.
In distributed multi-satellite scenarios, studies have also explored simultaneous fault estimation and cluster consensus control~\cite{10171456}.
A preliminary version of this work was presented in~\cite{xie2025fault}, where we proposed a streaming federated learning framework addressing communication constraints and data heterogeneity in satellite fault diagnosis. 
Those initial results demonstrated the feasibility of using dual-domain knowledge distillation and partial forgetting compensation mechanisms. 